% The talk's dot-plot grammar (spec, "Dot plots"). Input once in the preamble, after vnslides.
%
% \vndotplot[<keys>]{<dots.csv>}{<bars.csv>} draws one panel:
%   - one row per ticker picker, grouped by basket size, at y = -(row + \vndotGap group_index): 1.4 rows more
%     between groups, so each group's label has its own space, 1.1 rows above its first row, which leaves it clearly
%     more room below than the rows have between them, and more room again above it, from the group before;
%   - one dot per weighting rule, on the row's line: hedge on hollow, drawn first, and hedge off filled, drawn on top,
%     so where two dots meet the filled one stays whole (user, 2026-10-30);
%   - a pale bar from each row's lowest to highest dot (one per hedge setting; the two overlap on the row's line);
%   - optional vertical lines: SPY buy-and-hold (gold, dashed), equal-weight universe (grey, dotted) and a midline
%     (grey, dashed), each drawn only when its key is set;
%   - an optional label on the SPY or equal-weight line: at the top of the plot, just left of the line, running up
%     along it.
% Keys: xmin, xmax, xstep, xprecision (decimals on the tick labels), spy, ew, midline, spylabel and ewlabel (the
% lines' labels), width, height (of the plot area), rowlabels (true or false), labelsize (of the row labels, pt),
% xlabel and title (one line above the plot, in a box of fixed height and depth, so titles with and without
% descenders share a baseline). Panels set side by side line up on their x axes. Extra axis options for one panel go
% in the style vndot extra, which \vndotplot clears when it is done.
% The dots file has columns row, group_index, hedge and value; the bars file has row, group_index, label, group,
% hedge, lo and hi. Both come from scripts/30_talk_data.jl.
%
% \vndotreadrefs{<refs.csv>} sets \vndotSPY and \vndotEW from a refs file (columns name, value); a name that is not
% in the file leaves its macro empty.
% \vndotreadaxis{<axes.csv>}{<set>} sets \vndotAxisLo and \vndotAxisHi from the row of an axes file (columns set, lo,
% hi) whose set is <set>.
% \vndotlegend[<keys>] draws the legend on one line. Keys: legend spy and legend ew (true or false), which items to
% show; legend size, the font size in pt (8 under a plot, 7.5 on the footer line, to match the footer).
% \vndotlegendstacked draws the legend in four lines (hedge off and hedge on share the first) in a plot's upper right
% corner. A panel takes it through its extra axis options:
%   \pgfplotsset{vndot extra/.style={extra description/.code={\vndotlegendstacked}}}
% \vndothedgekey draws hedge on above hedge off, just right of a plot's top right corner, for a slide whose lines are
% labeled in the plot. A panel takes it the same way.
% \vndotupper{<text>} is <text> in capitals with γ left as it is: γ is a symbol on these slides, and capitals
% would turn it into Γ.

% the dots' radii and the range bars' width; a slide can make them smaller, inside a group around its frame
\def\vndotOffRadius{1.8pt}
\def\vndotOnRadius{1.7pt}
\def\vndotBarWidth{2.2pt}

% the colors (vndotoff, vndoton, vndotbar, vndotew, vndotgrid, vndotaxis, vndottick) come with the theme, from the
% deck palette in vnslides.sty

% the line and mark styles, shared by the plot and the legend
\tikzset{
  vndot spy/.style={draw=vnfigaccent, line width=1.1pt, dash pattern=on 4pt off 2.5pt},
  vndot ew/.style={draw=vndotew, line width=1pt, dotted},
  vndot mid/.style={draw=vndotew, line width=0.8pt, dash pattern=on 3pt off 2pt},
  vndot off/.style={circle, fill=vndotoff, draw=none, inner sep=0pt, minimum size=2*\vndotOffRadius},
  vndot on/.style={circle, fill=vnbg, draw=vndoton, line width=0.7pt, inner sep=0pt, minimum size=2*\vndotOnRadius},
  vndot line label/.style={rotate=90, anchor=south east, inner xsep=2pt, inner ysep=0pt, fill=vnbg, overlay,
    font=\fontsize{7.5}{9}\selectfont, text=vnslidebody},
}

% the space between a line's label and the line's center
\def\vndotLabelGap{3pt}

% the extra space between groups, in rows, and how far above its first row a group's label sits
\def\vndotGap{1.4}
\def\vndotGroupLift{1.1}

% extra axis options for the next panel only; \vndotplot clears this style when it is done
\pgfplotsset{vndot extra/.style={}}

% keep only the table rows whose column #1 reads #2 (the hedge column reads false or true)
\pgfplotsset{
  vndot keep/.style 2 args={
    x filter/.append code={%
      \edef\vndotCellA{\thisrow{#1}}%
      \edef\vndotCellB{#2}%
      \ifx\vndotCellA\vndotCellB
      \else
        \def\pgfmathresult{inf}%
      \fi
    },
  },
}

% the keys and their defaults; .estore in expands a value such as \vndotSPY when it is stored
\pgfkeys{
  /vndot/.is family, /vndot,
  xmin/.estore in=\vndotXmin,
  xmax/.estore in=\vndotXmax,
  xstep/.estore in=\vndotXstep,
  xprecision/.estore in=\vndotXprecision,
  spy/.estore in=\vndotSpyAt,
  ew/.estore in=\vndotEwAt,
  midline/.estore in=\vndotMidAt,
  spylabel/.store in=\vndotSpyLabel,
  ewlabel/.store in=\vndotEwLabel,
  width/.estore in=\vndotWidth,
  height/.estore in=\vndotHeight,
  rowlabels/.estore in=\vndotRowLabels,
  labelsize/.estore in=\vndotLabelSize,
  xlabel/.store in=\vndotXlabel,
  title/.store in=\vndotTitle,
  legend spy/.estore in=\vndotLegendSpy,
  legend ew/.estore in=\vndotLegendEw,
  legend size/.estore in=\vndotLegendSize,
}

\newcommand{\vndotdefaults}{%
  \pgfkeys{/vndot, xmin=0.95, xmax=1.75, xstep=0.1, xprecision=1, spy={}, ew={}, midline={}, spylabel={}, ewlabel={},
    width=8cm, height=4cm, rowlabels=true, labelsize=7.5, xlabel={}, title={}}%
}

\def\vndotTrue{true}
\def\vndotFalse{false}

% \vndotreadrefs{<refs.csv>}: \vndotSPY and \vndotEW from the refs file, looked up by name
\def\vndotNameSpy{SPY buy-and-hold}
\def\vndotNameEw{equal-weight universe}
\newcommand{\vndotreadrefs}[1]{%
  \def\vndotSPY{}%
  \def\vndotEW{}%
  \pgfplotstableread[col sep=comma]{#1}\vndotRefs
  \pgfplotstablegetrowsof{\vndotRefs}%
  \pgfmathtruncatemacro\vndotRefLast{\pgfplotsretval - 1}%
  \pgfplotsforeachungrouped \vndotI in {0,...,\vndotRefLast}{%
    \pgfplotstablegetelem{\vndotI}{name}\of\vndotRefs
    \edef\vndotRefName{\pgfplotsretval}%
    \pgfplotstablegetelem{\vndotI}{value}\of\vndotRefs
    \ifx\vndotRefName\vndotNameSpy
      \edef\vndotSPY{\pgfplotsretval}%
    \fi
    \ifx\vndotRefName\vndotNameEw
      \edef\vndotEW{\pgfplotsretval}%
    \fi
  }%
}

% \vndotreadaxis{<axes.csv>}{<set>}: \vndotAxisLo and \vndotAxisHi from the row of that set
\newcommand{\vndotreadaxis}[2]{%
  \def\vndotAxisLo{}%
  \def\vndotAxisHi{}%
  \edef\vndotWantSet{#2}%
  \pgfplotstableread[col sep=comma]{#1}\vndotAxes
  \pgfplotstablegetrowsof{\vndotAxes}%
  \pgfmathtruncatemacro\vndotAxisLast{\pgfplotsretval - 1}%
  \pgfplotsforeachungrouped \vndotI in {0,...,\vndotAxisLast}{%
    \pgfplotstablegetelem{\vndotI}{set}\of\vndotAxes
    \edef\vndotThisSet{\pgfplotsretval}%
    \ifx\vndotThisSet\vndotWantSet
      \pgfplotstablegetelem{\vndotI}{lo}\of\vndotAxes
      \edef\vndotAxisLo{\pgfplotsretval}%
      \pgfplotstablegetelem{\vndotI}{hi}\of\vndotAxes
      \edef\vndotAxisHi{\pgfplotsretval}%
    \fi
  }%
}

% \vndotupper{<text>}: capitals, with γ kept lower case
\ExplSyntaxOn
\cs_generate_variant:Nn \text_uppercase:n { V }
\cs_new_protected:Npn \vndotupper #1
  {
    \tl_set:Ne \l_tmpa_tl {#1}
    \tl_replace_all:Nnn \l_tmpa_tl { γ } { \NoCaseChange { γ } }
    \text_uppercase:V \l_tmpa_tl
  }
\ExplSyntaxOff

% \vndotplot[<keys>]{<dots.csv>}{<bars.csv>}
\newcommand{\vndotplot}[3][]{%
  \vndotdefaults
  \pgfkeys{/vndot, #1}%
  % the bar table drives the y range and the labels
  \pgfplotstableread[col sep=comma]{#3}\vndotBars
  \pgfplotstablegetrowsof{\vndotBars}%
  \pgfmathtruncatemacro\vndotLast{\pgfplotsretval - 1}%
  \pgfplotstablegetelem{\vndotLast}{row}\of\vndotBars
  \let\vndotLastRow\pgfplotsretval
  \pgfplotstablegetelem{\vndotLast}{group_index}\of\vndotBars
  \let\vndotLastGroup\pgfplotsretval
  \pgfmathsetmacro\vndotYmin{-(\vndotLastRow + \vndotGap*\vndotLastGroup) - 0.5}%
  % the top of the plot is the first group's label
  \pgfmathsetmacro\vndotYmax{-1 + \vndotGroupLift}%
  % a title only when one is set: an empty title still takes a line
  \ifx\vndotTitle\empty
    \pgfplotsset{vndot title/.style={}}%
  \else
    \pgfplotsset{vndot title/.style={title={\vndotTitle}}}%
  \fi
  \begin{tikzpicture}[baseline=(current axis.south)]
    \begin{axis}[
      width=\vndotWidth, height=\vndotHeight, scale only axis,
      clip=false,
      xmin=\vndotXmin, xmax=\vndotXmax, xtick distance=\vndotXstep,
      ymin=\vndotYmin, ymax=\vndotYmax,
      axis y line=none, axis x line*=bottom,
      x axis line style={draw=vndotaxis, -},
      xmajorgrids, major grid style={draw=vndotgrid, line width=0.5pt},
      major x tick style={draw=vndotaxis},
      tick label style={font=\fontsize{8}{9.6}\selectfont, color=vndottick},
      xticklabel style={/pgf/number format/assume math mode=true, /pgf/number format/fixed,
        /pgf/number format/fixed zerofill, /pgf/number format/precision=\vndotXprecision},
      xlabel={\vndotXlabel},
      vndot title,
      title style={font=\fontsize{8.5}{10}\selectfont, text=vnslideink, yshift=-2pt, text height=1.6ex,
        text depth=0.45ex},
      xlabel style={font=\fontsize{8.5}{10}\selectfont, color=vnslidebody, yshift=3pt},
      every axis plot/.append style={line join=round},
      vndot extra,
    ]
      % reference lines, behind the bars and dots
      \ifx\vndotMidAt\empty
      \else
        \draw[vndot mid] (axis cs:\vndotMidAt,\vndotYmin) -- (axis cs:\vndotMidAt,\vndotYmax);
      \fi
      \ifx\vndotEwAt\empty
      \else
        \draw[vndot ew] (axis cs:\vndotEwAt,\vndotYmin) -- (axis cs:\vndotEwAt,\vndotYmax);
      \fi
      \ifx\vndotSpyAt\empty
      \else
        \draw[vndot spy] (axis cs:\vndotSpyAt,\vndotYmin) -- (axis cs:\vndotSpyAt,\vndotYmax);
      \fi
      % a line's label runs up along the line's left side, \vndotLabelGap from the line's center, and its text ends at
      % the top of the plot; its field-colored box runs 2 pt past the text at each end, so a grid line under the label
      % stops short of the letters. It is an overlay, so it may reach past the plot's left edge without moving the panel
      \ifx\vndotEwLabel\empty
      \else
        \node[vndot line label] at ([xshift=-\vndotLabelGap, yshift=2pt]axis cs:\vndotEwAt,\vndotYmax)
          {\vndotEwLabel};
      \fi
      \ifx\vndotSpyLabel\empty
      \else
        \node[vndot line label] at ([xshift=-\vndotLabelGap, yshift=2pt]axis cs:\vndotSpyAt,\vndotYmax)
          {\vndotSpyLabel};
      \fi

      % bars: lowest to highest dot, drawn as an x error bar around the midpoint
      \addplot[draw=none, mark=none, vndot keep={hedge}{false},
        error bars/.cd, x dir=both, x explicit, error mark=none,
        error bar style={draw=vndotbar, line width=\vndotBarWidth, line cap=round}]
        table[col sep=comma, x expr={(\thisrow{lo} + \thisrow{hi})/2},
          y expr={-(\thisrow{row} + \vndotGap*\thisrow{group_index})},
          x error expr={(\thisrow{hi} - \thisrow{lo})/2}] {#3};
      \addplot[draw=none, mark=none, vndot keep={hedge}{true},
        error bars/.cd, x dir=both, x explicit, error mark=none,
        error bar style={draw=vndotbar, line width=\vndotBarWidth, line cap=round}]
        table[col sep=comma, x expr={(\thisrow{lo} + \thisrow{hi})/2},
          y expr={-(\thisrow{row} + \vndotGap*\thisrow{group_index})},
          x error expr={(\thisrow{hi} - \thisrow{lo})/2}] {#3};

      % dots: hedge on hollow first (filled with the field, so the bar does not show through), then hedge off filled on
      % top
      \addplot[only marks, mark=*, mark size=\vndotOnRadius, mark options={fill=vnbg, draw=vndoton, line width=0.7pt},
        vndot keep={hedge}{true}]
        table[col sep=comma, x=value, y expr={-(\thisrow{row} + \vndotGap*\thisrow{group_index})}] {#2};
      \addplot[only marks, mark=*, mark size=\vndotOffRadius, color=vndotoff,
        mark options={fill=vndotoff, line width=0pt}, vndot keep={hedge}{false}]
        table[col sep=comma, x=value, y expr={-(\thisrow{row} + \vndotGap*\thisrow{group_index})}] {#2};

      % row labels at each row, and each group's label above its first row
      \ifx\vndotRowLabels\vndotTrue
        \def\vndotPrevGroup{-1}%
        \pgfplotsforeachungrouped \vndotI in {0,...,\vndotLast}{%
          \pgfplotstablegetelem{\vndotI}{hedge}\of\vndotBars
          \edef\vndotHedge{\pgfplotsretval}%
          \ifx\vndotHedge\vndotFalse
            \pgfplotstablegetelem{\vndotI}{row}\of\vndotBars
            \let\vndotRow\pgfplotsretval
            \pgfplotstablegetelem{\vndotI}{group_index}\of\vndotBars
            \let\vndotGroupIndex\pgfplotsretval
            \pgfplotstablegetelem{\vndotI}{label}\of\vndotBars
            \let\vndotLabel\pgfplotsretval
            \pgfplotstablegetelem{\vndotI}{group}\of\vndotBars
            \let\vndotGroup\pgfplotsretval
            \pgfmathsetmacro\vndotY{-(\vndotRow + \vndotGap*\vndotGroupIndex)}%
            \edef\vndotCmd{\noexpand\node[anchor=east, inner sep=0pt, xshift=-6pt,
              font=\noexpand\fontsize{\vndotLabelSize}{\vndotLabelSize}\noexpand\selectfont, text=vnslideink]
              at (axis cs:\vndotXmin,\vndotY) {\vndotLabel};}%
            \vndotCmd
            % a new group: its label sits \vndotGroupLift row above this row
            \ifnum\vndotGroupIndex=\vndotPrevGroup
            \else
              \pgfmathsetmacro\vndotYgroup{\vndotY + \vndotGroupLift}%
              \edef\vndotCmd{\noexpand\node[anchor=east, inner sep=0pt, xshift=-6pt,
                font=\noexpand\fontsize{6}{7.2}\noexpand\selectfont\noexpand\addfontfeatures{LetterSpace=10},
                text=vnslidelink]
                at (axis cs:\vndotXmin,\vndotYgroup) {\noexpand\vndotupper{\vndotGroup}};}%
              \vndotCmd
              \edef\vndotPrevGroup{\vndotGroupIndex}%
            \fi
          \fi
        }%
      \fi
    \end{axis}
  \end{tikzpicture}%
  \pgfplotsset{vndot extra/.style={}}%
}

% \vndotlegend[<keys>]: the legend on one line
\newcommand{\vndotlegend}[1][]{%
  \pgfkeys{/vndot, legend spy=true, legend ew=true, legend size=8, #1}%
  \pgfmathsetmacro\vndotLegendSkip{1.2*\vndotLegendSize}%
  {\fontsize{\vndotLegendSize}{\vndotLegendSkip}\selectfont\color{vnslidebody}%
    \tikz[baseline=-0.6ex]\node[vndot off] {};~hedge off\hspace{1.2em}%
    \tikz[baseline=-0.6ex]\node[vndot on] {};~hedge on\hspace{1.2em}%
    \ifx\vndotLegendSpy\vndotTrue
      \tikz[baseline=-0.6ex]\draw[vndot spy] (0,0) -- (0.5,0);~SPY buy-and-hold\hspace{1.2em}%
    \fi
    \ifx\vndotLegendEw\vndotTrue
      \tikz[baseline=-0.6ex]\draw[vndot ew] (0,0) -- (0.5,0);~equal-weight universe\hspace{1.2em}%
    \fi
    one dot per weighting rule\par}%
}

% \vndotlegendstacked: the legend in four lines, in the plot's upper right. It is drawn in the axis's own picture (a
% nested \tikz would take on the axis scaling), on a field-colored box so the grid lines stop at it. (1,1) is the
% plot's top right corner, in axis description coordinates; the marks sit in a column at the left, the labels in a
% column 0.15 cm to their right, and each line is \vndotPitch below the one before.
\def\vndotPitch{9.6pt}
\newcommand{\vndotlegendstacked}{%
  \coordinate (vndotL) at (1,1);
  % the box behind the legend
  \fill[vnbg] ([xshift=-3.95cm, yshift=2pt]vndotL) rectangle ([yshift=-4*\vndotPitch - 2pt]vndotL);
  \begin{scope}[font=\fontsize{8}{9.6}\selectfont, every node/.append style={text=vnslidebody}]
    % line 1: hedge off and hedge on, each mark close to its own label
    \node[vndot off] at ([xshift=-3.41cm, yshift=-0.5*\vndotPitch]vndotL) {};
    \node[anchor=base west, inner sep=0pt] at ([xshift=-3.2cm, yshift=-0.5*\vndotPitch - 0.6ex]vndotL) {hedge off};
    \node[vndot on] at ([xshift=-1.55cm, yshift=-0.5*\vndotPitch]vndotL) {};
    \node[anchor=base west, inner sep=0pt] at ([xshift=-1.34cm, yshift=-0.5*\vndotPitch - 0.6ex]vndotL) {hedge on};
    % lines 2 and 3: the reference lines
    \draw[vndot spy] ([xshift=-3.85cm, yshift=-1.5*\vndotPitch]vndotL)
      -- ([xshift=-3.35cm, yshift=-1.5*\vndotPitch]vndotL);
    \node[anchor=base west, inner sep=0pt] at ([xshift=-3.2cm, yshift=-1.5*\vndotPitch - 0.6ex]vndotL)
      {SPY buy-and-hold};
    \draw[vndot ew] ([xshift=-3.85cm, yshift=-2.5*\vndotPitch]vndotL)
      -- ([xshift=-3.35cm, yshift=-2.5*\vndotPitch]vndotL);
    \node[anchor=base west, inner sep=0pt] at ([xshift=-3.2cm, yshift=-2.5*\vndotPitch - 0.6ex]vndotL)
      {equal-weight universe};
    % line 4: what a dot is
    \node[anchor=base west, inner sep=0pt] at ([xshift=-3.85cm, yshift=-3.5*\vndotPitch - 0.6ex]vndotL)
      {one dot per weighting rule};
  \end{scope}%
}

% \vndothedgekey: hedge on above hedge off, just right of the plot. The labels hang from the top of the plot, the
% second \vndotPitch below the first, 9.5 pt in from the plot's right edge; each mark sits 4.5 pt left of its label, at
% the label's mid height. It is an overlay, so the panel keeps its width.
\newcommand{\vndothedgekey}{%
  \coordinate (vndotK) at (1,1);
  \begin{scope}[overlay, font=\fontsize{7.5}{9}\selectfont, every node/.append style={text=vnslidebody}]
    \node[anchor=north west, inner sep=0pt] (vndotKon) at ([xshift=9.5pt]vndotK) {hedge on};
    \node[anchor=north west, inner sep=0pt] (vndotKoff) at ([yshift=-\vndotPitch]vndotKon.north west) {hedge off};
    \node[vndot on] at ([xshift=-4.5pt]vndotKon.mid west) {};
    \node[vndot off] at ([xshift=-4.5pt]vndotKoff.mid west) {};
  \end{scope}%
}
